
\documentclass{article}
\usepackage{colortbl}
\usepackage{makecell}
\usepackage{multirow}
\usepackage{supertabular}

\begin{document}

\newcounter{utterance}

\centering \large Interaction Transcript for game `codenames', experiment `risk\_high', episode 9 with reflexion{-}llama{-}3.1{-}8b{-}instruct.
\vspace{24pt}

{ \footnotesize  \setcounter{utterance}{1}
\setlength{\tabcolsep}{0pt}
\begin{supertabular}{c@{$\;$}|p{.15\linewidth}@{}p{.15\linewidth}p{.15\linewidth}p{.15\linewidth}p{.15\linewidth}p{.15\linewidth}}
    \# & $\;$A & \multicolumn{4}{c}{Game Master} & $\;\:$B\\
    \hline

    \theutterance \stepcounter{utterance}  
    & & & \multicolumn{2}{p{0.3\linewidth}}{
        \cellcolor[rgb]{0.9,0.9,0.9}{
            \makecell[{{p{\linewidth}}}]{
                \texttt{\tiny{[GM$|$GM]}}
                {'hidden': {'team': ['round', 'time', 'ketchup', 'circle', 'back'], 'innocent': ['eagle', 'olympus', 'lion'], 'opponent': ['worm', 'tooth', 'torch', 'ground', 'watch'], 'assassin': ['suit', 'paper', 'missile', 'pipe']}, 'revealed': {'team': {'team': [], 'innocent': [], 'opponent': [], 'assassin': []}, 'opponent': {'team': [], 'innocent': [], 'opponent': [], 'assassin': []}}}
            }
        }
    }
    & & \\ \\

    \theutterance \stepcounter{utterance}  
    & & \multicolumn{4}{p{0.6\linewidth}}{
        \cellcolor[rgb]{0.9,0.9,0.9}{
            \makecell[{{p{\linewidth}}}]{
                \texttt{\tiny{[P1$\langle$GM]}}
                \texttt{You have attempted this same task before and did not succeed. These are the plans you wrote after those attempts.} \\
\texttt{Trial \#1: New plan: In my last attempt, I provided a clue that was semantically related to the target words, but I failed to follow the format correctly. The Game Master rejected my message because it did not start with 'CLUE: ' followed by a new line and then 'TARGETS: ' followed by a comma{-}separated list of target words. Next time, I will make sure to follow the format correctly and provide a well{-}formed and legal response. For example, I will start with 'CLUE: shape' on a new line, followed by 'TARGETS: round, circle, trunk'.} \\
\\ 
\texttt{Trial \#2: New plan: In my last attempt, I provided a clue that was semantically related to the target words, but I failed to consider the opponent's team words and the distractor words. I should have chosen a clue that is more specific and less likely to be associated with the opponent's team words or the distractor words. Next time, I will analyze the lists more carefully and choose a clue that is more unique and related to the target words, but not to the opponent's team words or the distractor words. For example, I will choose a clue that is related to the shape or form of the target words, such as 'shape' or 'form'.} \\
\\ 
\texttt{Trial \#3: New plan: In my last attempt, I provided a clue that was semantically related to the target words, but I failed to consider the kill words. I should have chosen a clue that is not associated with the kill words, as providing a clue related to a kill word could end the game immediately. Next time, I will analyze the lists more carefully and choose a clue that is more unique and related to the target words, but not to the kill words. For example, I will choose a clue that is related to the shape or form of the target words, such as 'shape' or 'form', but not to words like 'suit' or 'pipe'.} \\
\\ 
\\ 
\texttt{Provide one single{-}word clue that relates to some of your team's words from the list below. You can choose to how many and to which words your clue relates to.} \\
\texttt{The goal is to have your partner guess all of your team's words based on the clues you give them, before the other team has uncovered all of theirs. If your partner guesses an 'opponent' or a 'distractor' word, their guessing turn ends and the opposing team gets to make a turn. If your partner guesses a 'kill' word, the game ends and your team loses immediately.} \\
\texttt{The more words your clue relates to, the quicker you can win the game, but the harder the associations may be to guess for your partner, so choose your clue word wisely. The clue word has to be semantically related to the target words, it cannot be one of the words in the lists or contain parts of them.} \\
\\ 
\texttt{Always give your single{-}word clue and your comma{-}separated list of related target words in the following format and make your answers as short as possible, never include any other text than is required in this form:} \\
\\ 
\texttt{CLUE: <WORD>} \\
\texttt{TARGETS: <TARGETS>} \\
\\ 
\texttt{Your team words are: round, time, ketchup, circle, back.} \\
\texttt{Your opponent's team words are: worm, tooth, torch, ground, watch.} \\
\texttt{Distractor words are: eagle, olympus, lion.} \\
\texttt{The kill words are: suit, paper, missile, pipe.} \\
            }
        }
    }
    & \\ \\

    \theutterance \stepcounter{utterance}  
    & \multicolumn{4}{p{0.6\linewidth}}{
        \cellcolor[rgb]{0.9,0.9,0.9}{
            \makecell[{{p{\linewidth}}}]{
                \texttt{\tiny{[P1$\rangle$GM]}}
                \texttt{CLUE: round} \\
\texttt{TARGETS: round, circle} \\
            }
        }
    }
    & & \\ \\

    \theutterance \stepcounter{utterance}  
    & & & \multicolumn{2}{p{0.3\linewidth}}{
        \cellcolor[rgb]{0.9,0.9,0.9}{
            \makecell[{{p{\linewidth}}}]{
                \texttt{\tiny{[GM$|$GM]}}
                {'player': 'Cluegiver', 'type': 'clue is morphologically related to word on the board', 'utterance': 'CLUE: round\nTARGETS: round, circle', 'message': "Your clue 'round' is morphologically similar to the word round, please choose another clue word.", 'clue': 'round', 'similar_board_word': 'round'}
            }
        }
    }
    & & \\ \\

    \theutterance \stepcounter{utterance}  
    & \multicolumn{4}{p{0.6\linewidth}}{
        \cellcolor[rgb]{0.95,0.95,0.95}{
            \makecell[{{p{\linewidth}}}]{
                \texttt{\tiny{[P2$\rightarrow$P1]}}
                \texttt{Post{-}episode Reflexion (not sent to GM):} \\
\texttt{New plan: In my last attempt, I provided a clue that was semantically related to the target words, but I failed to consider the kill words and the opponent's team words simultaneously. I should have chosen a clue that is more unique and related to the target words, but not to the opponent's team words, the distractor words, or the kill words. Next time, I will analyze the lists more carefully and choose a clue that is more specific and less likely to be associated with the opponent's team words, the distractor words, or the kill words. For example, I will choose a clue that is related to the shape or form of the target words, such as 'shape' or 'form', but not to words like 'suit', 'paper', 'missile', or 'pipe'.} \\
            }
        }
    }
    & & \\ \\

\end{supertabular}
}

\end{document}
